\section{Experimenteller Aufbau}
\label{sec:experimental-setup}

\subsection{Inhaltliche Stichpunkte}
\begin{itemize}
  \item \textbf{Netzwerk und Knoten}
  \begin{itemize}
    \item Acht simulierte Knoten: Node A bis Node H.
    \item In der dokumentierten Aggregation werden je Zielknoten sieben
          Peer-Berichte ausgewertet.
    \item Die genaue Graphstruktur und ob reale Netzwerkkommunikation
          simuliert wird, vor der finalen Beschreibung festlegen.
  \end{itemize}
  \item \textbf{Datensatz und Aufteilung}
  \begin{itemize}
    \item Je Knoten 100\,000 Zeilen in 100 Blöcken; je Block 70\,\% Training
          und 30\,\% Test.
    \item Insgesamt: 560\,000 Trainings- und 240\,000 Testzeilen über alle
          Knoten.
    \item Klassenverteilung im Training: none 305\,900; syn-flood 149\,100;
          Backdoor 105\,000. Im Test: none 131\,100; syn-flood 63\,900;
          Backdoor 45\,000.
    \item Die Datensätze enthalten pro Klasse Original-, TabGAN- und
          DAE-Beispiele; die im fertigen Datensatz beobachteten Anteile
          stehen in Kapitel~\ref{sec:synthetic-data}.
    \item Aufteilung laut Notebook~02 blockweise mit 70/30-Split;
          genaue Reihenfolge und Zufallsbehandlung der Zeilen dokumentieren.
  \end{itemize}
  \item \textbf{Merkmale und Klassen}
  \begin{itemize}
    \item Numerische Merkmale: shunt\_voltage, bus\_voltage\_V,
          current\_mA und power\_mW; zusätzlich State.
    \item Zielklassen: Backdoor, none und syn-flood.
  \end{itemize}
  \item \textbf{Varianten und Metriken}
  \begin{itemize}
    \item Primärer Notebook~07-Vergleich auf 240\,000 Testfällen:
          Baseline ohne Kontext gegenüber Model~2 mit Full Context + State.
    \item Baseline: Accuracy 0{,}930625, Macro-F1 0{,}911805,
          Weighted-F1 0{,}928974, Log Loss 0{,}182478 und
          ROC-AUC (Macro OVR) 0{,}986353.
    \item Model~2: Accuracy 0{,}962204, Macro-F1 0{,}953685,
          Weighted-F1 0{,}961924, Log Loss 0{,}098996 und
          ROC-AUC (Macro OVR) 0{,}994932.
    \item Zusätzlich Precision, Recall und F1 je Klasse sowie
          Konfusionsmatrizen und ROC-Kurven auswerten.
  \end{itemize}
\end{itemize}

\subsection{Noch zu erklären}
\begin{itemize}
  \item Notebook~02 teilt innerhalb 1000-Zeilen-Blöcken auf; erläutern,
        ob zeitlich benachbarte Zeilen in Training und Test landen können
        und wie sich dies auf die Unabhängigkeit der Testdaten auswirkt.
  \item Die 240\,000 Testzeilen stammen aus den finalen Knoten-Datensätzen
        mit gemischten Datenquellen. Von dem separaten Original-Testdatensatz
        aus Notebook~01b unterscheiden und die jeweilige Verwendung erklären.
  \item Die 100\,000 Zeilen je Knoten und die Klassenverteilungen mit den
        erzeugten Szenarien und Originaldaten abgleichen.
  \item Hardware, Softwareversionen, RF-Parameter, Zufalls-Seeds und
        konkrete Normalisierungsschritte nachvollziehbar angeben.
  \item Precision, Recall, F1 und ROC-AUC sind in den Ergebnisnotebooks
        ausgewiesen; eine Falsch-Positiv-Rate und ein Signifikanztest sind
        dort nicht direkt ausgewiesen und müssen gegebenenfalls ergänzt werden.
\end{itemize}

\subsection{Grafikbedarf}
\begin{itemize}
  \item \textbf{Vorhanden:} \url{evaluation/plots/test_attack_sequence_nodes_a_h.png}
        zur zeitlichen Einordnung der Testblöcke; Trainingsplot nur bei
        zusätzlichem Erkenntniswert verwenden.
  \item \textbf{Erstellen:} Netzwerkdiagramm mit Node A bis H und der
        tatsächlich genutzten Peer-Verbindung.
  \item \textbf{Erstellen:} übersichtliche Tabelle oder Grafik mit
        Klassenverteilung und Datenquellen im Trainings- und Testsatz.
\end{itemize}
